% ================================================================
% APS / REVTeX 4.2 RESEARCH PAPER TEMPLATE
% Based on the supplied two-column APS paper structure.
% Replace bracketed prompts with your own writing and results.
% Delete instructional comments before submitting.
% ================================================================
\documentclass[
 reprint,
 amsmath,amssymb,
 aps,
]{revtex4-2}

\usepackage{graphicx} % Figures / external images
\usepackage{dcolumn}  % Decimal-aligned table columns
\usepackage{tikz}     % Optional diagrams
\usepackage{xcolor}
\usepackage{float}
\usepackage{comment}
\usepackage[colorlinks=true,linkcolor=black,citecolor=black,urlcolor=black]{hyperref}
\usepackage{bm}

\begin{document}

\title{[Concise paper title: system studied, key variable, or principal finding]}
\author{[Author 1, Author 2, Author 3]}
\affiliation{[Department, institution, city, postcode, country]}
\date{\today}

\begin{abstract}
% Write this last, ideally as one self-contained paragraph.
% Include: (1) research context / problem; (2) precise aim;
% (3) experimental or computational approach; (4) main quantitative
% result, with uncertainty where appropriate; (5) key conclusion.
% Avoid unexplained abbreviations, lengthy background and citations.
[Briefly state the question, what you did, the principal numerical
finding, and what the finding means.]
\end{abstract}
\maketitle

\section{Introduction}
\label{sec:introduction}

% Paragraph 1: motivate the physical phenomenon and its relevance.
[Introduce the physical problem and explain why it matters.]

% Paragraph 2: provide only the background needed to understand your
% investigation. Define the model, physical quantities and assumptions.
% Cite suitable primary literature or review articles, e.g.:
% \cite{example_article} (replace this example key with a real source).
[Summarise the relevant model, prior research and essential context.]

% Paragraph 3: derive / present the theoretical prediction being tested.
% Define EVERY symbol, assumptions and expected dependence.
[Explain the theory underlying the expected trend, and identify the
prediction against which you will compare your data.]

\begin{equation}
\label{eq:theory}
 y = A x^{\alpha},
\end{equation}
% Replace the illustrative equation with the theory for your system.
% Explain the meanings of y, x, A and alpha, or delete the example.

% Final paragraph: state the gap, research question, hypotheses and
% what the paper will test, rather than just listing techniques.
[State the specific research question and the quantitative tests used
to answer it.]

\section{Methodology}
\label{sec:methodology}

% Describe the experimental apparatus OR computational model.
% State parameters, units, boundary / initial conditions, sampling
% methods, controls, number of trials and how reproducibility is ensured.
[Describe the setup and procedure in enough detail to reproduce it.]

% Explain observable definitions and criteria used to classify outcomes.
% Distinguish operational definitions from ideal theoretical quantities.
[Define the quantities measured and the rules used to classify or
extract observations.]

% Describe data processing / parameter estimation / fitting and how
% uncertainties were estimated and propagated. Mention any relevant
% independence, convergence or finite-size checks.
[Describe the analysis workflow and statistical treatment.]

% OPTIONAL SCHEMATIC: replace the placeholder with a real diagram,
% or remove this figure entirely. A real image could instead use:
% \includegraphics[width=\columnwidth]{figures/setup.pdf}
\begin{figure}[htbp]
    \centering
    \fbox{\parbox[c][3cm][c]{0.88\columnwidth}{\centering
    [Insert experimental schematic, simulation geometry, or workflow]}}
    \caption{[Describe what is shown; define symbols and explain
    the relevant physical or computational setup.]}
    \label{fig:setup}
\end{figure}

\section{Results \& Discussion}
\label{sec:results}

% Present important results in a logical sequence, usually with the
% primary figure/table first. Report values with units and uncertainty.
% Focus prose on what the data show, not on restating every plotted point.
[Present the main measured trend and its estimated uncertainty.]

% REPLACE this example with a real data plot. Keep figure references
% in the prose, e.g. Fig.~\ref{fig:main_result}.
\begin{figure}[htbp]
    \centering
    \fbox{\parbox[c][3.4cm][c]{0.88\columnwidth}{\centering
    [Insert main result plot with axes, units, error bars and legend]}}
    \caption{[Describe the plotted quantities, conditions, and error
    bars. Define abbreviations so the plot can stand alone.]}
    \label{fig:main_result}
\end{figure}

% Example table layout. Replace with data, or delete if unnecessary.
\begin{table}[htbp]
\caption{[Comparison of fitted parameters, conditions or key results;
state uncertainty conventions and units.]}
\label{tab:results}
\begin{ruledtabular}
\begin{tabular}{lcc}
Condition & Measured quantity & Prediction \\
\hline
{[Case A]} & [value $\pm$ uncertainty] & [value] \\
{[Case B]} & [value $\pm$ uncertainty] & [value] \\
\end{tabular}
\end{ruledtabular}
\end{table}

% Compare directly against the prediction given in the Introduction.
% Where possible, quantify agreement/disagreement: fitted exponents,
% residuals, confidence intervals, goodness-of-fit or other apt tests.
[Quantify agreement with theory, including numerical differences and
whether those differences are resolved by the uncertainty.]

% Interpret the physical mechanism. Compare to appropriate literature.
% Distinguish credible systematic limitations from speculative issues.
% Discuss refinements only when they meaningfully explain observations.
[Interpret the results physically, compare with published work and
discuss substantive limitations or possible model corrections.]

\section{Conclusion}
\label{sec:conclusion}

% Answer the question from the end of the Introduction directly.
% State the principal quantitative result, whether the predicted
% relation was supported, and its physical interpretation.
% Mention the most important limitation / implication if warranted.
% Do NOT introduce new evidence, broad unrelated applications,
% extensive literature review or unsupported claims.
[In a short paragraph, answer the research question, give the main
numerical finding with uncertainty, state the strength of agreement
with theory, and identify its principal implication.]

% Keep bibliography.bib beside this file. Replace its example records
% with real sources and cite their BibTeX keys using \cite{key}.
\bibliography{bibliography}

\end{document}
